\begin{figure}[!htbp]
\centering
\resizebox{0.98\textwidth}{!}{%
\begin{tikzpicture}[
  risk/.style={draw=gray!35, fill=white, rounded corners=7pt,
    line width=0.62pt, align=center, font=\scriptsize,
    text width=2.45cm, minimum height=1.12cm, inner sep=5pt},
  site/.style={risk, draw=orange!88!black, fill=orange!7},
  path/.style={risk, draw=orange!80!black, fill=orange!9},
  agent/.style={risk, draw=accent!70, fill=blue!5},
  review/.style={risk, draw=accent!55, fill=blue!3},
  note/.style={draw=gray!35, fill=gray!3, rounded corners=5pt,
    line width=0.50pt, font=\scriptsize\bfseries, align=center,
    text width=10.4cm, inner xsep=5pt, inner ysep=4pt, text=black!70}
]

\node[site] (r1) at (0,1.35) {source site down};
\node[path] (r2) at (2.90,1.35) {hosting page removed};
\node[path] (r3) at (5.80,1.35) {embedded page inaccessible};
\node[path] (r4) at (8.70,1.35) {server/player unavailable};

\node[agent] (r5) at (0,-0.10) {anti-bot drift};
\node[review] (r6) at (2.90,-0.10) {CDN/provider lookup gaps};
\node[agent] (r7) at (5.80,-0.10) {partial-run telemetry};
\node[review] (r8) at (8.70,-0.10) {human/legal review};

\node[note] at (4.35,-1.50)
  {Review action: rerun time-bound cases, inspect trace and screenshot evidence, then
  decide whether the limitation is external, agentic, or legal-review related.};

\end{tikzpicture}%
}
\caption{Operational risk map for interpreting evaluation results.}\label{fig:limitations-risk-map}
\end{figure}
